\documentclass[8pt]{beamer}	%声明文档类型为beamer幻灯片
\mode <presentation>

%Antibes Bergen Berkeley Berlin Boadilla Copenhagen Darmstadt Dresden Frankfurt Goettingen Hannover Ilmenau Juanlespins Madrid Malmoe Marburg Montpellier Paloalto Pittsburgh Rochester Singapore

%◦ 没有导航条: default, AnnArbor, Bergen, Boadilla, CambridgeUS, . . .
%◦ 带树形导航条: Antibes, JuanLesPins, Montpellier
%◦ 带侧边导航条: Berkeley, Goettingen, Hannover, Marburg, PaloAlto
%◦ 带微型导航条: Berlin, Ilmenau, Darmstadt, Dresden, Frankfurt, . . .
%◦ 带节小节标题: Copenhagen, Luebeck, Malmoe, Warsaw
\usetheme{Madrid}%%Warsaw [secheader]
%\usecolortheme{beaver}
\beamersetaveragebackground{black!10}

\setbeamercovered{transparent}
%\setbeamertemplate{navigation symbols}{}

\usefonttheme{professionalfonts}
\useinnertheme{circles}%{rectangles}
\setbeamertemplate{itemize item}{$\circledast$}%{\checkmark}


\setbeamertemplate{title page}
{
%  \vbox{}
  \vfill
  \begin{centering}
  	\vskip0.6em\par%
  	\includegraphics[height=2.0cm]{ustc_logo_fig_new.eps}
    \vskip0.4em\par%
    \begin{beamercolorbox}[sep=8pt,center,shadow=true,rounded=true]{title}
      	\usebeamerfont{title}\inserttitle\par%
      	\ifx\insertsubtitle\@empty%
      	\else%
        \vskip0.25em%
        {\usebeamerfont{subtitle}\usebeamercolor[fg]{subtitle}\insertsubtitle\par}%
      	\fi%
    \end{beamercolorbox}%
    \vskip1em\par
    \begin{beamercolorbox}[sep=8pt,center]{author}
      \usebeamerfont{author}
 %     \insertauthor
      {
        \begin{tabular}{cc}
  			报告人：&\insertauthor\\
  			导\quad 师：& \advisor
  		\end{tabular}
      }
    \end{beamercolorbox}
    \begin{beamercolorbox}[sep=10pt,center]{institute}
      \usebeamerfont{institute}\insertinstitute
    \end{beamercolorbox}
    \begin{beamercolorbox}[sep=10pt,center]{date}
      \usebeamerfont{date}\insertdate
    \end{beamercolorbox}\vskip1.5em
  %  {\usebeamercolor[fg]{titlegraphic}\inserttitlegraphic\par}
  	\end{centering}
  	\vfill
}

%\hypersetup{pdfpagemode=FullScreen} % makes your presentation go automatically to full screen

\usepackage{fontspec,xunicode,xltxtra}
\usepackage{xeCJK} %设定让幻灯片支持中文

\usepackage{multimedia}	%让文档支持多媒体
\usepackage{graphics}	%让文档支持图片
\usepackage{hyperref}	%让文档支持超链接
\usepackage{booktabs}	%让文档支持三线表格
\usepackage{amsmath}	%ams可以让文档支持数学公式
\usepackage{amsfonts}
\usepackage{amssymb}
\usepackage{color}
\usepackage{graphicx,psfrag}
\usepackage{epsfig}

\setCJKmainfont[BoldFont={SimSun}, ItalicFont={SimSun}]{SimSun}	%设置缺省中文字体，根据自己系统里的字体选择合适字体
\setCJKsansfont{WenQuanYi Micro Hei}	%设置中文无衬线字体
\setCJKmonofont{WenQuanYi Micro Hei}	%设置中文等宽字体

\setCJKfamilyfont{song}{WenQuanYi Micro Hei}%自定义字体，可以用来在幻灯片中改变字体
\setCJKfamilyfont{hei}{SimHei}	%自定义字体，可以用来在幻灯片中改变字体
\setCJKfamilyfont{kai}{SimKai}
\setCJKfamilyfont{fs}{WenQuanYi Micro Hei}

%\setCJKmainfont[BoldFont=SimHei,ItalicFont=KaiTi_GB2312]{KaiTi_GB2312}
%\setCJKsansfont{SimHei}	%设置中文无衬线字体
%\setCJKmonofont{FangSong_GB2312}
%
%\setCJKfamilyfont{song}{NSimSun}
%\setCJKfamilyfont{hei}{SimHei}
%\setCJKfamilyfont{kai}{KaiTi_GB2312}   % XP对应 KaiTi_GB2312，Vista对应KaiTi，注意根据系统切换
%\setCJKfamilyfont{fs}{FangSong_GB2312} % XP对应 FangSong_GB2312，Vista对应FangSong，注意根据系统切换


\newcommand{\song}{\CJKfamily{song}}    % 宋体
\newcommand{\fs}{\CJKfamily{fs}}        % 仿宋体
\newcommand{\kai}{\CJKfamily{kai}}      % 楷体
\newcommand{\hei}{\CJKfamily{hei}}      % 黑体

\setmainfont[Mapping=tex-text]{Times New Roman}
\setsansfont[Mapping=tex-text]{Arial}%[Mapping=tex-text]{Arial}%[BoldFont={Courier New Bold}]{Courier New}
\setmonofont[Mapping=tex-text]{Courier New}%[BoldFont={Arial:style=Bold}]{Arial}


\title{EAST密度极限破裂数据库的建立及其预测}	%幻灯片标题
\author[陈俊杰]{陈俊杰}
\def\advisor{秦宏}

\institute[中国科学技术大学]{{\normalsize 工程与应用物理系 \\ 中国科学技术大学，物理学院}}%作者单位
%\date{\today}
\date{2019.09.12}

\setcounter{tocdepth}{1}

\begin{document}
\begin{frame}[plain]
\titlepage
\end{frame}

\AtBeginSection[]{
  \frame{
    \frametitle{大纲}
    \tableofcontents[currentsection]
  }
}

\begin{frame}
\frametitle{内容提要}
\tableofcontents
\end{frame}

%%%%%%%%%%%%%%%%%%%%%%%%%%%%%%%%%%%%%%%%%%%%%%%%%%%%%%%%%%%%%%%%%%%%%%%%%%%%%%%

\section{概述}
\begin{frame}
\frametitle{概述}
研究意义
\end{frame}

\begin{frame}
\frametitle{概述}
研究内容与方法
\end{frame}
%%%%%%%%%%%%%%%%%%%%%%%%%%%%%%%%%%%%%%%%%%%%%%%%%%%%%%%%%%%%%%%%%%%%%%%%%%%%%%%

\section{存在的问题}

\begin{frame}
\frametitle{一二三四五问题}
\begin{itemize}
\item 一二三四五
\item \alert{六七八九十}
\end{itemize}
\pause
一二三四五

\end{frame}

%%%%%%%%%%%%%%%%%%%%%%%%%%%%%%%%%%%%%%%%%%%%%%%%%%%%%%%%%%%%%%%%%%%%%%%%%%%%%%%

\section{方法一}

\begin{frame}
\frametitle{USTC}

\begin{itemize}
\item<1-> USTC,ABCD EFG HIJK LMN OPQ RST UVW XYZ
\item<2-> USTC,abcd efg hijk lmn opq rst uvw xyz
\end{itemize}


\end{frame}

%%%%%%%%%%%%%%%%%%%%%%%%%%%%%%%%%%%%%%%%%%%%%%%%%%%%%%%%%%%%%%%%%%%%%%%%%%%%%%%

\section{方法二}

\begin{frame}
\frametitle{公式}
$$M=m_{ubm}+Dy$$
\end{frame}

%%%%%%%%%%%%%%%%%%%%%%%%%%%%%%%%%%%%%%%%%%%%%%%%%%%%%%%%%%%%%%%%%%%%%%%%%%%%%%%

\section{方法三}

\begin{frame}
\frametitle{Block}
\begin{block}{Block1}
	\begin{itemize}
		\item 方法三
        \item 方法三
	\end{itemize}
\end{block}
\pause
\begin{block}{Block1}
	\begin{itemize}
		\item 方法三
        \item 方法三
	\end{itemize}
\end{block}
\end{frame}

%%%%%%%%%%%%%%%%%%%%%%%%%%%%%%%%%%%%%%%%%%%%%%%%%%%%%%%%%%%%%%%%%%%%%%%%%%%%%%%

\section{总结与展望}

\begin{frame}
\frametitle{实验结果}
\begin{columns}
	\begin{column}{0.30\textwidth}
		\epsfig{figure=ustc_logo_fig_new,width=\textwidth}
	\end{column}
	\begin{column}{0.50\textwidth}
		\begin{block}{结论}
			\begin{itemize}
			\item 实验结果如左图所示
			\item 实验结果如左图所示
			\item 实验结果如左图所示
			\item 实验结果如左图所示
			\item 实验结果如左图所示
			\item 实验结果如左图所示
			\end{itemize}
		\end{block}
	\end{column}
\end{columns}
\end{frame}


%%%%%%%%%%%%%%%%%%%%%%%%%%%%%%%%%%%%%%%%%%%%%%%%%%%%%%%%%%%%%%%%%%%%%%%%%%%%%%%


\begin{frame}
\centerline{\Huge 谢谢!}
\end{frame}

\end{document} 